\documentclass[11pt,a4paper]{article}

\usepackage[utf8]{inputenc}
\usepackage[T1]{fontenc}
\usepackage[french]{babel}
\usepackage[margin=2cm]{geometry}
\usepackage{graphicx}
\usepackage{hyperref}
\usepackage{enumitem}
\usepackage{xcolor}

\hypersetup{colorlinks=true, linkcolor=black, urlcolor=blue!50!black}
\setlist{nosep}

\title{\vspace{-1.5cm}Jeu d'\'Echec Yalta --- Rapport technique}
\author{Gaël Amoros}
\date{Avril 2026}

\begin{document}
\maketitle
\vspace{-0.8cm}

\section{Contexte et analyse du problème}

Le Yalta est une variante d'échecs à trois joueurs jouée sur un plateau hexagonal de
96~cases, chaque joueur disposant des 16~pièces classiques (soit 48~pièces au total).
Le projet consiste à concevoir et implémenter ce jeu en C++ en appliquant une
modélisation orientée objet rigoureuse et trois design patterns, dont le
\emph{Modèle-Vue-Contrôleur} imposé par le sujet. Une intelligence artificielle
fondée sur min-max / alpha-beta multi-threadée est attendue dans une itération
ultérieure et sort donc du périmètre de ce rapport intermédiaire.

Les difficultés principales identifiées en phase d'analyse sont~: (i)~la
représentation des 96~cases du plateau Yalta, qui n'est pas une grille cartésienne
classique et impose une coordonnée à trois composantes (\emph{secteur},
\emph{colonne}, \emph{rang})~; (ii)~les règles de déplacement des pièces lorsqu'elles
franchissent la frontière entre deux secteurs, en particulier pour les fous et les
reines dont les diagonales « tournent »~; (iii)~la gestion d'une partie à trois
joueurs, avec l'élimination d'un joueur maté et la poursuite de la partie entre les
deux restants.

\section{Choix de conception}

\subsection{Architecture MVC}

L'application est structurée en trois couches étanches~:
\begin{itemize}
    \item \textbf{Modèle}~: \texttt{Plateau}, \texttt{Partie}, \texttt{Piece} et ses
    sous-classes, \texttt{Coup}. Ne dépend d'aucune bibliothèque graphique, ce qui
    le rend testable unitairement et réutilisable pour brancher l'IA.
    \item \textbf{Vue}~: \texttt{FenetrePrincipale}, \texttt{VuePlateau} (un
    \texttt{QGraphicsScene}), \texttt{VuePiece}. Se contente de rendre l'état et
    d'émettre des signaux Qt lorsqu'une case est cliquée.
    \item \textbf{Contrôleur}~: \texttt{ControleurJeu} reçoit les signaux de la Vue,
    valide les coups auprès du Modèle et propage les mises à jour.
\end{itemize}

Le pattern est renforcé par le mécanisme \emph{signals/slots} de Qt, qui joue le
rôle d'un Observer intégré~: toute modification du modèle émet un signal que la vue
écoute pour se redessiner, sans couplage direct.

\subsection{Strategy pour les déplacements}

Chaque pièce délègue le calcul de ses coups possibles à une implémentation de
l'interface \texttt{IStrategieDeplacement} (\texttt{DeplacementPion},
\texttt{DeplacementTour}, etc.). Cela évite un \texttt{switch} géant sur le type de
pièce dans le moteur, et --- surtout --- préparera l'intégration future de
stratégies de décision IA (\texttt{MinMax}, \texttt{AlphaBeta}) qui remplaceront
\texttt{ControleurJeu} comme source de coups pour un joueur donné sans modification
du modèle.

\subsection{Factory pour la création des pièces}

\texttt{PieceFactory} centralise deux responsabilités~: instancier une pièce à
partir de son type et de sa couleur (avec injection de la bonne Strategy), et
produire la disposition initiale des 48~pièces sur les trois camps Yalta. Cette
centralisation évite la duplication de la logique de câblage et facilite la
gestion de la \emph{promotion du pion} qui, concrètement, demande de remplacer une
pièce par une autre au même emplacement.

\section{Diagrammes}

Trois diagrammes UML accompagnent cette analyse (dossier
\texttt{docs/diagrammes/})~:
\begin{itemize}
    \item \textbf{Diagramme de classes} --- structure statique complète avec les
    trois patterns mis en évidence~;
    \item \textbf{Diagramme d'états-transitions} --- cycle d'un tour de jeu, avec
    les états terminaux spécifiques à la variante 3~joueurs~;
    \item \textbf{Diagramme de cas d'utilisation} --- acteurs et interactions
    attendues, avec l'acteur \emph{Joueur IA} déjà identifié pour l'itération
    suivante.
\end{itemize}

% \section{Fonctionnalités implémentées}
% \section{Difficultés rencontrées}
% \section{Améliorations futures (IA min-max multi-threadée)}
% --- à compléter en fin de projet ---

\end{document}
